\documentclass[11pt]{article}
\usepackage[margin=1in]{geometry}
\usepackage{amsmath,amssymb,amsthm,mathtools}
\usepackage{array,booktabs}
\usepackage[hidelinks]{hyperref}

\newtheorem{theorem}{Theorem}[section]
\newtheorem{lemma}[theorem]{Lemma}
\newtheorem{proposition}[theorem]{Proposition}
\newcommand{\spE}{s^+}
\newcommand{\smE}{s^-}
\newcommand{\sig}{\sigma}
\newcommand{\one}{\mathbf 1}

\title{Positive Square Energy of Every Decacyclic Cactus}
\author{}
\date{26 July 2026}

\begin{document}
\maketitle

\begin{abstract}
Let $G$ be a connected cactus of cyclomatic rank ten on $n$ vertices.  We
prove $s^+(G)>n$.  The sharp block-additive doubly nonnegative estimate reduces
the problem exactly to $T^9Q$ and $T^8PP$, where $T=C_3$, $P=C_5$, and $Q$ is
an arbitrary cycle.  Exact colored cluster-partition censuses give
$96=92+4$ and $180=170+10$ proper rows.  Exhaustive reduced-tree pruning and marked finite
certificates reduce their structural rows to $A_9\mid Q$, $T^8P\mid P$, and
$P\mid A_8\mid P$.  The first has $3624=3621+3$ marked-entry classes; the
second has $11689=11586+100+3$; and the third has
$11689=11674+15$.  In the one-cluster case the full $T^9Q$ census has at most
$8049$ color-preserving incidences and an exact three-row hostile frontier.
The $T^8PP$ census is $30386=30377+9$, and all nine exceptions have explicit
positive replacements.  Separations use actual bridges or consecutive cycle
intervals, and every connector remnant, incidence branch, shared cut, and
attached tree receives exactly one final owner.  A tempting packing-one guard
is false on the new three-cut kernels; we instead open a leaf cycle and apply
packing one only to the retained common-hub packet.  No open two-pivot winding
assertion is used.
\end{abstract}

\section{Definitions, inputs, and firewall}

All graphs are finite, simple, and undirected.  A cactus is a connected graph
whose blocks are bridges or cycles; its cyclomatic rank is the number of cyclic
blocks.  For adjacency eigenvalues $\lambda_1,\ldots,\lambda_n$, set
\[
 \spE(G)=\sum_{\lambda_i>0}\lambda_i^2,
 \qquad \smE(G)=\sum_{\lambda_i<0}\lambda_i^2,
 \qquad \sig(G)=\spE(G)-|V(G)|.
\]
For every induced partition $V(G)=V_1\mathbin{\dot\cup}\cdots
\mathbin{\dot\cup}V_r$, positive square energy is superadditive:
\begin{equation}\label{eq:super}
 \spE(G)\geq\sum_i\spE(G[V_i]),
 \qquad \sig(G)\geq\sum_i\sig(G[V_i]).
\end{equation}
This is the theorem of Akbari, Kumar, Mohar, and Pragada~\cite{AKMP}.
Since a connected rank-ten graph has $|E|=n+9$, the theorem below proves the
strict positive-square-energy conjecture for this class.

Write $T=C_3$, $P=C_5$, and
\[
 \delta=\sec(\pi/5)-1=\sqrt5-2<\frac14,
 \qquad \delta_q=\sec(\pi/q)-1<1\quad(q\geq4).
\]
A cycle $Q=C_q$ is \emph{hostile} when $q\equiv1\pmod4$.  We use the proved
rank-nine cactus theorem, including arbitrary attached trees~\cite{RankNine},
the common-cut and one-hostile-cycle packet theorems~\cite{CommonCut,Packing},
and the exact rank-ten artifacts listed in Appendix~\ref{app:reproduction}.

The proof has a strict firewall.  It does not use the proposed two-pivot phase
or winding assertion of~\cite{TwoPivot}; the exact two-terminal reduction there
does not prove the needed integrated argument bound.  It does not infer
packing number one from a nearest-cycle or Voronoi territory.  It does not
contract a bridge connector analytically, split a retained shared cut between
owners, use edge monotonicity, or replace a cyclic argument by a principal
$\arctan$.  Every invocation of packing one below displays the actual common
hub of the retained triangles.  The two new apparent two-pivot rows are closed
by opening a leaf cycle at exact tree cost $-1$.

\section{Packets, cycle order, and ownership}

The rank-nine theorem and its packet infrastructure give, uniformly over
arbitrary finite trees attached at arbitrary packet vertices,
\begin{align}
 \sig(H)&\geq0 &&\text{for connected ranks two and three},\label{eq:low0}\\
 \sig(H)&>0 &&\text{for connected ranks four through nine},\label{eq:lowpos}\\
 \sig(P)&\geq-\delta,& \sig(Q)&\geq-\delta_q,\label{eq:hostile}\\
 \sig(TP)&>1-\delta,& \sig(TQ)&>0,& \sig(TTQ)&\geq0,\label{eq:mixed}\\
 \sig(PP)&>0,& \sig(TPP)&>\frac32.\label{eq:pp}
\end{align}
For a connected shared-cut cluster $A_r$ of $r$ triangles we use
\begin{equation}\label{eq:triangle}
 \sig(A_r)>b_r,
 \qquad (b_1,\ldots,b_8)=(0,1,2,3,2,1,0,0).
\end{equation}

\begin{lemma}[One-hostile-cycle packing one]\label{lem:packing}
Let the cyclic blocks of a connected cactus $H$ be one hostile $Q=C_q$ and
$a\geq1$ triangles, no two of which are vertex-disjoint.  The $Q$ may meet the
triangular part at its root or be joined to it by an arbitrary positive path,
and arbitrary trees may occur everywhere.  Then
\[
 \spE(H)-\smE(H)>-2\delta_q,
 \qquad \sig(H)>a-\delta_q.
\]
\end{lemma}

\begin{proof}
Eliminate every tree off the cyclic spine by its nonnegative matching message.
In the grouped Sachs expansion on $it$, write
$\Psi_H(t)=R+2i(B-A)$, where $B$ is the matching partition after deleting
$Q$, and $A$ is the sum after additionally deleting one triangle.  Packing
one excludes every term containing two triangles.  Splitting matchings by use
of the $Q$ boundary gives
\[
 R-Z_q(t)(B-A)=E+LB+Z_q(t)A+
 4\sum_{T\cap Q=\varnothing}Z_{H-(V(Q)\cup V(T))}>0.
\]
Thus the continuous argument is strictly smaller than the isolated hostile
cycle argument.  The signed Coulson identity gives the first inequality.
Since $|E(H)|=|V(H)|+a$ and $\spE+\smE=2|E|$, the second follows.  This is the
standalone argument of~\cite{Packing}, not a two-pivot estimate.
\end{proof}

\begin{lemma}[Common-cut packets]\label{lem:common}
For a one-vertex bouquet, with arbitrary attached trees,
\begin{align}
 \sig(T^kQ)&>k-\delta_q &&(q\equiv1\pmod4),\label{eq:ccq}\\
 \sig(T^kQ)&>k &&(q\text{ even or }q\equiv3\pmod4),\label{eq:ccgood}\\
 \sig(T^kP)&>k-\delta,\label{eq:ccp}\\
 \sig(T^kPP)&>k+1-\frac4{3\sqrt{13}}.\label{eq:ccpp}
\end{align}
\end{lemma}

\begin{proof}
After exact tree elimination, the normalized determinant is the scalar sum
\[
 W(t)=a_x+\sum_j\frac{B_j+\omega_{\ell_j}}{A_j},
 \qquad \omega_\ell=-2i^{-\ell},
\]
with $A_j,B_j>0$.  Comparing its real and imaginary parts with the isolated
hostile lobe gives~\eqref{eq:ccq}--\eqref{eq:ccp}.  For two pentagons the exact
coefficient inequality
\[
 2R\geq t(t^4+7t^2+9)(A_1+A_2)
\]
and $\int_0^\infty4/(t^4+7t^2+9)\,dt=2\pi/(3\sqrt{13})$ give
\eqref{eq:ccpp}.  The complete derivation and its 1290-term exact coefficient
check are~\cite{CommonCut,Coefficient}.
\end{proof}

Join cyclic blocks when they share a cut vertex; the components are
\emph{shared-cut clusters}.  In the block-cut tree contract each cluster, take
the minimal subtree spanning the resulting marked cluster nodes, and suppress
only unmarked degree-two nodes.  Components off this cyclic hull are trees with
a unique hull attachment and stay with that attachment's owner.  The resulting
\emph{reduced cluster tree} has only marked leaves.  A reduced edge records a
nonempty chain of actual bridges; it is never treated as an analytic edge.

\begin{lemma}[Cycle order and interval realization]\label{lem:cycleorder}
Let a cycle $C$ receive one, two, or three owner marks from components of the
graph outside $C$.  Their cyclic order determines a partition of $C$ into
nonempty proper consecutive vertex intervals, with coincident marks retained
together.  Assigning each open arc and either boundary vertex according to
that order gives disjoint induced territories.  A later router operation may
refine one existing interval but never crosses an earlier boundary.
\end{lemma}

\begin{proof}
Delete the edges of $C$ and inspect the components rooted at its vertices.
Each has one root, so its owner is forced.  For two distinct marks, cut one
edge in each complementary arc; for three marks, cut one edge in each of the
three arcs.  Move each cut within its open arc until every interval is
nonempty.  The cactus property excludes any route between two such components
outside $C$, so the resulting vertex sets are induced.  Refinement repeats
the same construction inside one interval and therefore preserves all earlier
boundaries.
\end{proof}

\begin{lemma}[Global ownership]\label{lem:ownership}
Suppose connected subtrees of the reduced cluster tree are assigned to packet
owners and every selected router is split as in Lemma~\ref{lem:cycleorder}.
Then the assignments lift to an induced, disjoint, exhaustive vertex
partition of the original cactus.  Every bridge remnant, connector entry,
shared cut, incidence branch, and off-hull tree has exactly one final owner.
If several router operations are sequential, replacing each provisional owner
when its interval is refined gives the same unique final ownership.
\end{lemma}

\begin{proof}
Cut only actual bridges between reduced-tree territories and use the unique
path property to assign every connector vertex.  Inside a cluster apply
Lemma~\ref{lem:cycleorder}; a branch at a cut follows that cut, while a tree at
a private cycle vertex follows its interval.  These rules cover every vertex
and agree at every boundary.  Induction on refinements proves the final-owner
claim.
\end{proof}

\section{The exact rank-ten frontier}

For an edge-containing graph put
\[
 \kappa(G)=\sup_{M\succeq0,\ M\geq0,\ \one^TM\one>0}
 \frac{4(\sum_{uv\in E(G)}\sqrt{M_{uv}})^2}{\one^TM\one}.
\]
The sharp cactus DNN theorem~\cite{DNN,LTZ} gives
\[
 \kappa(G)=b+\sum_i\kappa(C_{\ell_i}),
 \qquad
 \kappa(C_\ell)=\begin{cases}\ell,&\ell\text{ even},\\
 \ell\sec^2(\pi/(2\ell)),&\ell\text{ odd},\end{cases}
 \qquad \smE(G)\leq\kappa(G).
\]
Set $\epsilon_\ell=0$ for even $\ell$ and
$\epsilon_\ell=\ell\tan^2(\pi/(2\ell))$ for odd $\ell$.  Block counting gives
\begin{equation}\label{eq:dnn}
 \sig(G)\geq9-\sum_{i=1}^{10}\epsilon_{\ell_i}.
\end{equation}
The odd sequence decreases, $\epsilon_3=1$, and with
$a=\epsilon_5=5-2\sqrt5$ one has
$3a<2$, $2a>1$, and $\epsilon_5+\epsilon_7<1$.  If $k$ cycles are
nontriangles, a nonpositive margin forces $k(1-a)\leq1$, hence $k\leq2$.
For $k=2$, the last comparison leaves only two pentagons.  Therefore the exact
frontier is
\begin{equation}\label{eq:residuals}
 T^9Q\quad(q\geq3,\text{ including }Q=T),
 \qquad T^8PP.
\end{equation}
Both genuinely remain: their DNN margins are $-\epsilon_q$ and
$1-2\epsilon_5=4\sqrt5-9$.

The exact colored cluster census~\cite{Frontier} is
\begin{center}
\begin{tabular}{lrrrr}
\toprule
residual&all partitions&proper&direct&structural\\
\midrule
$T^9Q$&97&96&92&4\\
$T^8PP$&181&180&170&10\\
\bottomrule
\end{tabular}
\end{center}
Here a cluster color is its triangle count together with its number of
distinguished cycles.  Every direct row has a positive exact packet ledger.
The structural rows are
\begin{align*}
T^9Q:
&\quad Q|T|T|T|T|T|T|T|T|T,
\quad Q|T|T|T^7,
\quad Q|T|T^8,
\quad Q|T^9;\\
T^8PP:
&\quad P|P|T|T|T|T|T|T|T|T,
\quad P|P|T|T^7,
\quad P|P|T^8,\\
&\quad P|T|T|T|T|T|T|T^2P,
\quad P|T|T|T|T|T|T^3P,\\
&\quad P|T|T|T|T|T^4P,
\quad P|T|T|T|T^5P,
\quad P|T|T|T^6P,\\
&\quad P|T|T^7P,
\quad P|T^8P.
\end{align*}

\begin{lemma}[Exact pruning of the ten $T^8PP$ structural rows]
\label{lem:t8ppprune}
For every topology of the reduced cluster tree, the ten structural rows in
the $T^8PP$ census have the following exhaustive ordered reductions.  An
all-triangle leaf is used first; its triangle count is at most eight and its
complement has rank at least two.
\begin{center}
\small
\begin{tabular}{cll}
\toprule
&colored row&outcome if no all-triangle cluster is a leaf\\
\midrule
1&$P_0|P_1|8T$&$TP+$ connected rank eight\\
2&$P_0|P_1|T|A_7$&$TP+$ connected rank eight\\
3&$P_0|P_1|A_8$&$P_0|A_8|P_1$\\
4&$P_1|6T|T^2P_0$&$TP+$ connected rank eight\\
5&$P_1|5T|T^3P_0$&$TP+$ connected rank eight\\
6&$P_1|4T|T^4P_0$&$TP+$ connected rank eight\\
7&$P_1|3T|T^5P_0$&$TP+$ connected rank eight\\
8&$P_1|2T|T^6P_0$&$TP+$ connected rank eight\\
9&$P_1|T|T^7P_0$&$TP+$ connected rank eight\\
10&$P_1|T^8P_0$&$T^8P_0|P_1$\\
\bottomrule
\end{tabular}
\end{center}
Every displayed separation is made on actual bridges and lifts to disjoint,
induced, exhaustive territories with a unique owner for every connector
remnant and attached tree.
\end{lemma}

\begin{proof}
If an all-triangle cluster $A_r$ is a leaf, cut the first actual bridge on its
unique reduced-tree chain.  Equation~\eqref{eq:triangle} makes that territory
strict, and its connected complement has rank $10-r\in\{2,\ldots,9\}$, hence
is nonnegative (and is strict from rank four onward).

Assume there is no such leaf.  Every leaf then contains a pentagon.  Since
there are only two pentagons, the reduced tree has exactly two leaves and is a
path, with one pentagon at each end.  In rows 1, 2, and 4--9, one end is the
singleton $P_1$ and its adjacent marked cluster is a singleton triangle.
Their terminal segment is a connected $TP$ territory; cutting the first
actual bridge on its other side leaves a connected rank-eight cactus.  Both
territories are strict.  In row 3 the only remaining path is
$P_0|A_8|P_1$, with $A_8$ between the two pentagons.  In row 10 the unique
path is the marked last-bridge endpoint $T^8P_0|P_1$.

For completeness, reduced trees with $k$ marked clusters have only unmarked
Steiner vertices of degree at least three and therefore at most $k-2$ such
vertices.  Adding a marked leaf either at a vertex or by subdividing an edge
generates every reduced tree inductively.  The exact colored topology program
enumerates all $156263$ resulting classes in the ten rows and verifies that
$156253$ use the leaf alternative, eight use the terminal $TP$ alternative,
and one uses each marked endpoint~\cite{T8PPTopology}.  Thus no path, Steiner,
or high-degree disconnected topology is omitted.

On every bridge cut, each chain remnant follows its side, and every off-hull
tree follows its unique attachment.  Thus Lemma~\ref{lem:ownership} applies.
The ten rows are exhausted.  A detailed graph-level audit is~\cite{T8PPPrune}.
\end{proof}

\begin{proposition}[Global cluster synthesis]\label{prop:global}
Every proper partition in the two residual families is strict.  After direct
packet rows and reduced-tree pruning, the only finite endpoints requiring a
marked certificate are
\begin{equation}\label{eq:endpoints}
 A_9\mid Q,\qquad T^8P_0\mid P_1,
 \qquad P_0\mid A_8\mid P_1.
\end{equation}
\end{proposition}

\begin{proof}
For each of the first three $T^9Q$ structural rows, every reduced tree has a
triangular leaf, since all leaves are marked and $Q$ is the only nontriangular
cluster.  Cutting its first actual bridge gives a strict triangular territory;
the connected complement has rank at most nine and is nonnegative, with a
strict packet on one side in every case.  The fourth row is exactly
$A_9\mid Q$.  Thus all Steiner topologies are exhausted.  For $T^8PP$ apply
Lemma~\ref{lem:t8ppprune}; its row table proves every structural topology or
leaves exactly $T^8P\mid P$ and $P\mid A_8\mid P$.  Together with the 92 and
170 direct packet rows this exhausts the four and ten frozen structural rows.
The marked closures of the three endpoints are proved next.
\end{proof}

\section{The marked disconnected endpoints}

\subsection{The endpoint \texorpdfstring{$A_9\mid Q$}{A9 | Q}}

Project the actual connector to its first shared cut or private triangle port.
The exact marked-entry certificate~\cite{T9QMarked} gives
\[
 355\text{ unmarked }A_9\text{ incidences},\qquad
 3624=3621+3\text{ canonical marked classes},\qquad
 6745\text{ physical marked positions}.
\]
Every one of the 3621 direct classes has a legal triangle-router split with
the $TQ$, $TTQ$, strict rank-three-through-nine, and triangular packet ledgers.
If the entry is private, its interval remains with the real connector and $Q$.

Two residual marks are the common-cut nine-triangle bouquet, entered at the
hub or privately; Lemma~\ref{lem:common} or Lemma~\ref{lem:packing} gives
$\sig>9-\delta_q$.  In the third, a triangle router joins a seven-petal hub to
one leaf triangle and the entry is private on the router.  Keep the connector
with the router.  The retained eight triangles share the real hub and the
rooted $Q$, so Lemma~\ref{lem:packing} gives $\sig>8-\delta_q>0$.  Thus all
3624 classes close.  The executable freezes the three marked signatures and
their orbit multiplicities $1,1,18$, checks $6725+20=6745$ before quotienting,
and assigns every connector, cut, interval, and attachment a final owner.

\subsection{The last bridge \texorpdfstring{$T^8P\mid P$}{T8P | P}}

The exact certificate~\cite{T8PLast} reports
\[
 2392\text{ }T^8P\text{ incidences},\quad
 1105\text{ with }P_0\text{ an incidence leaf},\quad
 11689=11586+100+3.
\]
The 100 finite replacements use zero, one, two, and three routers in
$2,9,73,16$ classes.  Final-owner validation is part of the executable.  The
three residual marks lie on one two-cut kernel: a router triangle joins a
common-cut $A_7$ fan to $P_0$.  Depending on the entry, the repairs are
\[
 A_7+PP>0,
 \qquad P+\text{packing-one }A_7P>7-2\delta,
 \qquad \text{open }P_1+\text{packing-one }T^8P_0>7-\delta.
\]
For the opening, the four private vertices of $P_1$ and all trees rooted there
form a nonempty induced tree of surplus $-1$; the complement has the displayed
eight-triangle common hub.  All 11689 classes close.

\subsection{The two-interface endpoint
\texorpdfstring{$P\mid A_8\mid P$}{P | A8 | P}}

The two labelled connectors may enter at cuts or private vertices and may
coincide.  Exhaustive placement on 126 unmarked $A_8$ incidences gives 36414
ordered placements and
\[
 11689=11674+15
\]
canonical classes~\cite{A8Interface}.  The best plans use zero, one, two, and
three routers in $6,10844,838,1$ classes.  The ordinary ledger is
$c-e-2\delta$, where $c$ is the retained integer credit and $e$ counts naked
private-interface intervals; it accepts exactly the rows with $c-e\geq1$.

Nine residuals are labelled forms of one two-hub kernel: a router $R$ joins six
petals at $x$ and one leaf triangle at $y$.  Its five surgeries give
$A_6+T+PP$, $A_6+TP+P$, common-cut $T^6P+T+P$, or packing-one
$T^6P+T+P$, with margins $>1$, $>2-2\delta$, or $>6-2\delta$.
The other six are bouquet interface orbits.  They close by opening one remote
pentagon and retaining packing-one $A_8P$; splitting an entered triangle to
give $P+$ common-cut $T^7P$; or producing $A_7+PP$ or $A_6+P+P$.
The weakest margin is $1-2\delta>0$.  Lemmas~\ref{lem:cycleorder} and
\ref{lem:ownership} justify every labelled interval and connector owner.

This proves Proposition~\ref{prop:global}, hence every residual graph with more
than one shared-cut cluster.

\section{One fully shared cluster}

\subsection{The complete \texorpdfstring{$T^9Q$}{T9Q} census}

The color-preserving incidence census~\cite{T9QFull} is
\begin{center}
\small
\begin{tabular}{lrrrrrrrrr|r}
\toprule
$Q$ capacity&1&2&3&4&5&6&7&8&9&total\\
\midrule
3&1&12&91&406&1178&2115&2250&1246&275&7574\\
4&1&12&91&412&1203&2187&2361&1340&306&7913\\
5&1&12&91&412&1208&2201&2393&1372&321&8011\\
6&1&12&91&412&1208&2204&2400&1383&327&8038\\
7&1&12&91&412&1208&2204&2402&1386&330&8046\\
8&1&12&91&412&1208&2204&2402&1387&331&8048\\
$\geq9$&1&12&91&412&1208&2204&2402&1387&332&8049\\
\bottomrule
\end{tabular}
\end{center}
Every nonhostile incidence except the common bouquet has an ordinary legal
split.  Actual hostility means $q\equiv1\pmod4$: among the displayed finite
rows this includes $q=5$, while $q=7$ is nonhostile.  The certificate also runs
the weaker hostile ledger at capacity seven as a conservative check.  For an
actual hostile $q\geq9$, the incidence universe is the saturated capacity-nine
universe.  In each of these weaker-ledger audits the exact frontier has the
same three frozen signatures:
\begin{description}
\small
\item[Q1] \path{X(Q()T()T()T()T()T()T()T()T()T())}
\item[Q2] \path{T(X(Q())X(T()T()T()T()T()T()T()T()))}
\item[Q3] \path{T(X(Q())X(T())X(T()T()T()T()T()T()T()))}
\end{description}
Q1 is common-cut and satisfies $\sig>9-\delta_q$.  In Q2 all nine triangles
contain one hub and Lemma~\ref{lem:packing} gives the same bound.

Q3 is the essential firewall example.  The saturated router meets $Q$, one
leaf triangle, and a common-hub branch of seven triangles.  The nine triangles
do \emph{not} have packing number one, so a global packing-one or Voronoi guard
would be false.  Open the leaf triangle: its two private vertices, their edge,
and every tree rooted there form a nonempty induced tree $E$, so $\sig(E)=-1$.
The complement $H$ contains the router and seven other triangles, all sharing
the real hub, and $Q$ rooted at the router's other cut.  Therefore
\[
 \sig(G)\geq\sig(H)+\sig(E)>(8-\delta_q)-1
 =7-\delta_q>0.
\]
The executable checks all three signatures in each actual or conservatively
hostile regime, canonical
orbit sums, router degrees, the openable leaf, and the retained common hub.  It
materializes the two induced territories in Q3 and assigns every router cut and
cycle attachment a final owner.  Exact integer comparisons with
$0<\delta_q<1$ certify the three strict margins, under both normal Python and
\texttt{python -O}.  Thus every fully shared $T^9Q$ cactus is strict.

\subsection{The complete \texorpdfstring{$T^8PP$}{T8PP} census:
\texorpdfstring{$30386=30377+9$}{30386 = 30377 + 9}}

The exact census and ordinary split result are
\begin{center}
\begin{tabular}{lrrrrrrrrr|r}
\toprule
cuts&1&2&3&4&5&6&7&8&9&total\\
\midrule
all&1&19&204&1155&3990&8135&9615&5843&1424&30386\\
safe&0&17&200&1154&3989&8135&9615&5843&1424&30377\\
exceptions&1&2&4&1&1&0&0&0&0&9\\
\bottomrule
\end{tabular}
\end{center}
The nine frozen signatures and replacements are:
\begin{center}
\small
\begin{tabular}{cll}
\toprule
code&$c$&replacement and positive margin\\
\midrule
E1&1&common-cut $T^8PP$: $>9-4/(3\sqrt{13})$\\
E2&2&open leaf $P$; common-cut $T^8P$: $>7-\delta$\\
E3&2&$P+$ common-cut $T^7P$: $>7-2\delta$\\
E4&3&open leaf $P$ and leaf $T$; common-cut $T^7P$: $>5-\delta$\\
E5&3&open one $P$; packing-one $T^8P$: $>7-\delta$\\
E6&3&$P+T+$ common-cut $T^6P$: $>6-2\delta$\\
E7&3&$P+P+A_6$: $>1-2\delta$\\
E8&4&$P+P+T+A_5$: $>2-2\delta$\\
E9&5&$P+P+T+T+A_4$: $>3-2\delta$\\
\bottomrule
\end{tabular}
\end{center}
Their signatures, in certificate order, are
\begin{description}
\small
\item[E1] \path{X(P()P()T()T()T()T()T()T()T()T())}
\item[E2] \path{P(X(P())X(T()T()T()T()T()T()T()T()))}
\item[E3] \path{T(X(P())X(P()T()T()T()T()T()T()T()))}
\item[E4] \path{P(X(P())X(T())X(T()T()T()T()T()T()T()))}
\item[E5] \path{T(X(P())X(P())X(T()T()T()T()T()T()T()))}
\item[E6] \path{T(X(P())X(P()T()T()T()T()T()T())X(T()))}
\item[E7] \path{X(T()T()T()T()T()T()T(X(P()))T(X(P())))}
\item[E8] \path{X(T()T()T()T()T()T(X(P()))T(X(P())X(T())))}
\item[E9] \path{X(T()T()T()T()T(X(P())X(T()))T(X(P())X(T())))}
\end{description}
E5 is the second firewall row.  Its ordinary split gives only $P+P+A_7$,
whose lower ledger is negative.  Open one leaf pentagon at cost $-1$ instead.
The central router and seven hub triangles then form eight triangles with
packing number one, joined to the other pentagon at the router's other vertex;
Lemma~\ref{lem:packing} gives $>8-\delta$, hence $>7-\delta$ after opening.
In E4, open the leaf pentagon and the leaf triangle at their private pivots,
each at cost $-1$, leaving common-cut $T^7P$ and margin $>5-\delta$.
The other rows use the displayed actual common cuts and legal triangle-router
intervals.  The nine-row verifier~\cite{Nine} regenerates all 30386 incidences,
checks all signatures and edges, and certifies all nine exact ledgers.

\section{Main theorem}

\begin{theorem}\label{thm:main}
Every connected cactus $G$ of cyclomatic rank ten, on $n$ vertices, satisfies
\[
 \boxed{\spE(G)>n}.
\]
\end{theorem}

\begin{proof}
The sharp DNN inequality~\eqref{eq:dnn} proves every cycle multiset except
\eqref{eq:residuals}.  For a residual, the reduced cluster tree either has more
than one shared-cut cluster or one.  In the first case, the exact
$96=92+4$ and $180=170+10$ colored audits, Proposition~\ref{prop:global}, and
the three marked endpoint certificates prove strict positivity.  In the
second, the complete $T^9Q$ census and its three hostile closures, or the
$30386=30377+9$ $T^8PP$ census and nine replacements, prove it.

Every global cut is an actual bridge cut.  Every cycle split follows cyclic
order and uses nonempty proper consecutive intervals.  Sequential operations
refine one owner, and Lemma~\ref{lem:ownership} gives every connector remnant,
cut, branch, and attached tree one final owner.  Common-cut estimates are used
only at displayed common pivots.  Packing one is used only after direct
verification of the retained hub; the two false-guard rows are instead opened
at exact cost $-1$.  No open two-pivot statement enters.  Hence
$\sig(G)>0$, which is the assertion.
\end{proof}

\appendix
\section{Exact reproduction}\label{app:reproduction}

The seven rank-ten programs require Python 3.10 or newer and use the standard
library with exact integer or \texttt{Fraction} arithmetic.  From the
repository root run
\begin{verbatim}
python3 research/rank-ten-cactus-frontier-census.py
python3 -O research/rank-ten-cactus-frontier-census.py
python3 research/decacyclic-t9q-marked-entry-certificate.py
python3 -O research/decacyclic-t9q-marked-entry-certificate.py
python3 research/decacyclic-t9q-incidence-certificate.py
python3 -O research/decacyclic-t9q-incidence-certificate.py
python3 research/decacyclic-t8p-last-bridge-census.py
python3 -O research/decacyclic-t8p-last-bridge-census.py
python3 research/decacyclic-t8-two-interface-census.py
python3 -O research/decacyclic-t8-two-interface-census.py
python3 research/decacyclic-fully-shared-nine-exceptions.py
python3 -O research/decacyclic-fully-shared-nine-exceptions.py
python3 research/decacyclic-t8pp-reduced-tree-topology.py
python3 -O research/decacyclic-t8pp-reduced-tree-topology.py
\end{verbatim}
Each program is run normally and under optimization.  Its authoritative
acceptance checks are explicit and fail closed; the dedicated marked and
replacement programs check frozen totals, canonical classes or digests,
interval ownership, router degrees, exact ledgers, and closure.  Build this
manuscript with
\begin{verbatim}
bash scripts/build-paper.sh all-decacyclic-cacti
\end{verbatim}

\begin{thebibliography}{99}
\bibitem{AKMP} S.~Akbari, H.~Kumar, B.~Mohar, and S.~Pragada,
\emph{A linear lower bound for the square energy of graphs}, Electron. J.
Combin. 32(3) (2025), Paper P3.53.
\bibitem{LTZ} Y.~Liu, Q.~Tang, and S.~Zhang,
\emph{The positive and negative square-energy conjecture}, arXiv:2607.18031,
2026.
\bibitem{DNN} Companion manuscript, \emph{The Optimized Liu--Tang--Zhang
Constant for Cactus Graphs}, 2026; \path{sharp-cactus-dnn/paper.tex}.
\bibitem{RankNine} Companion manuscript, \emph{Positive Square Energy of Every
Nonacyclic Cactus}, 2026; \path{all-nonacyclic-cacti/paper.tex}.
\bibitem{Frontier} \emph{Exact rank-ten cactus frontier census}, 26 July 2026;
\path{research/rank-ten-cactus-frontier-census-2026-07-26.md} and
\path{research/rank-ten-cactus-frontier-census.py}.
\bibitem{T9QMarked} Exact marked certificate,
\path{research/decacyclic-t9q-marked-entry-certificate.py}.
\bibitem{T9QFull} Exact fully shared certificate and closure note,
\path{research/decacyclic-t9q-incidence-certificate.py} and
\path{research/decacyclic-t9q-proof-and-gap-2026-07-26.md}.
\bibitem{T8PLast} Exact last-bridge certificate,
\path{research/decacyclic-t8p-last-bridge-census.py}.
\bibitem{A8Interface} \emph{A marked $A_8$ interface theorem for the rank-nine
to rank-ten step}, 26 July 2026;
\path{research/rank-ten-marked-a8-interface-theorem-2026-07-26.md} and
\path{research/decacyclic-t8-two-interface-census.py}.
\bibitem{T8PPPrune} \emph{Graph-level pruning of the ten disconnected
$T^8PP$ structural rows}, 26 July 2026;
\path{research/decacyclic-t8pp-structural-pruning-lemma-2026-07-26.md}.
\bibitem{T8PPTopology} Exact reduced-tree topology census,
\path{research/decacyclic-t8pp-reduced-tree-topology.py} and
\path{research/decacyclic-t8pp-reduced-tree-topology-2026-07-26.md}.
\bibitem{Nine} Exact nine-exception verifier and census note,
\path{research/decacyclic-fully-shared-nine-exceptions.py} and
\path{research/decacyclic-rank-ten-router-census-2026-07-26.md}.
\bibitem{Packing} \emph{Rooted packing-one hostile-cycle lemma for octacyclic
provenance}, 26 July 2026;
\path{research/octacyclic-packing-one-hostile-cycle-lemma-2026-07-26.md}.
\bibitem{CommonCut} \emph{Common-cut residual bouquets: an exact rooted
Schur--Sachs inequality}, 26 July 2026;
\path{research/common-cut-bouquet-rooted-schur-2026-07-26.md}.
\bibitem{Coefficient} Exact coefficient verifier,
\path{positive-square-energy/experiments/c5_bouquet_matching_certificate.py}.
\bibitem{TwoPivot} \emph{Two-pivot Schur--Sachs triangular cactus reduction},
open winding boundary, 26 July 2026;
\path{research/two-pivot-schur-sachs-triangular-cactus-2026-07-26.md}.
\end{thebibliography}
\end{document}
